\chapter[High Availability y Replicación]{\emj{🔄}~High Availability y Replicación}

\section[Streaming Replication]{\emj{📡}~Streaming Replication}

PostgreSQL usa \textbf{WAL (Write-Ahead Log)} como fuente de verdad para replicación.
El primary envía registros WAL al standby en tiempo real (streaming), que los aplica para mantener una copia casi idéntica.

\begin{teoriabox}{Arquitectura de Streaming Replication}
\begin{itemize}
  \item \textbf{Primary:} acepta lecturas y escrituras. Genera WAL en \texttt{pg\_wal/}.
  \item \textbf{WAL Sender:} proceso en primary que envía WAL al standby por TCP.
  \item \textbf{WAL Receiver:} proceso en standby que recibe WAL y lo escribe localmente.
  \item \textbf{Startup Process:} en standby, aplica WAL continuamente (recovery loop).
  \item \textbf{Replication lag:} delay entre cambio en primary y aplicación en standby. Medido en bytes (write lag) y tiempo (replay lag).
\end{itemize}
\end{teoriabox}

La configuración de streaming replication requiere cambios en el primary y luego crear el standby desde un base backup.
El primary necesita permitir conexiones de replicación (en \texttt{pg\_hba.conf}) y tener el nivel de WAL correcto.
El standby se inicializa con \texttt{pg\_basebackup}, que copia el data directory completo y establece la conexión inicial.

\begin{lstlisting}[caption={Configurar streaming replication (Primary)}]
-- postgresql.conf en PRIMARY:
wal_level = replica            -- mínimo para replicación
max_wal_senders = 5            -- procesos WAL sender simultáneos
wal_keep_size = 1GB            -- retener WAL para standbys lentos
hot_standby = on               -- standbys pueden servir lecturas (default on)

-- pg_hba.conf en PRIMARY (agregar):
-- host replication replicator 10.0.0.2/32 scram-sha-256

-- Crear usuario de replicación:
CREATE USER replicator WITH REPLICATION ENCRYPTED PASSWORD 'securepass';
\end{lstlisting}

\begin{lstlisting}[caption={Configurar streaming replication (Standby)}]
-- Crear standby con pg_basebackup:
pg_basebackup -h primary_host -U replicator -D /var/lib/postgresql/data \
  --wal-method=stream --checkpoint=fast --progress

-- postgresql.conf en STANDBY:
hot_standby = on
hot_standby_feedback = on      -- evita que autovacuum en primary cancele queries en standby

-- postgresql.conf o recovery.conf (PG 12+: primary_conninfo en postgresql.conf):
primary_conninfo = 'host=primary_host port=5432 user=replicator password=securepass'

-- Crear archivo de señal para modo standby (PG 12+):
touch /var/lib/postgresql/data/standby.signal
\end{lstlisting}

El monitoreo continuo del lag es crítico en un clúster con replicación asíncrona.
Un standby con 10 minutos de lag no puede servir como failover sin pérdida de datos.
PostgreSQL expone el lag en bytes (diferencia de LSN) y en tiempo (cuánto tardaría en alcanzar al primary).
El lag en tiempo es más intuitivo para definir SLAs de RPO, pero el lag en bytes es más preciso para detectar el problema.

\begin{lstlisting}[caption={Monitorear replication lag}]
-- En el PRIMARY:
SELECT client_addr, state, sent_lsn, write_lsn, flush_lsn, replay_lsn,
       write_lag, flush_lag, replay_lag,
       sync_state
FROM pg_stat_replication;

-- replay_lag > 0: standby está por detrás
-- sync_state: 'async' (default), 'sync', 'quorum'

-- En el STANDBY:
SELECT now() - pg_last_xact_replay_timestamp() AS replication_delay;
SELECT pg_is_in_recovery();          -- true si es standby
SELECT pg_last_wal_replay_lsn();     -- último WAL aplicado
\end{lstlisting}

\section[Replicación Síncrona vs Asíncrona]{\emj{⚖️}~Replicación Síncrona vs Asíncrona}

\begin{teoriabox}{Modos de commit}
\begin{itemize}
  \item \textbf{Asíncrona (default):} primary confirma commit sin esperar al standby. RPO $>$ 0 (puede haber pérdida de datos). Latencia mínima.
  \item \textbf{Síncrona:} primary espera que el standby confirme haber recibido (y opcionalmente aplicado) el WAL antes de retornar el commit. RPO = 0 (sin pérdida). Latencia mayor.
  \item \textbf{synchronous\_commit = remote\_apply:} espera a que el standby \emph{aplique} el WAL. Garantía más fuerte pero latencia máxima.
\end{itemize}
\end{teoriabox}

La replicación síncrona garantiza que cuando el primary confirma un commit al cliente, ese commit ya está en al menos un standby.
El primary bloquea el commit hasta recibir la confirmación del standby designado.
Esto elimina el riesgo de pérdida de datos en failover (RPO = 0), pero añade latencia proporcional al round-trip time entre primary y standby: en un datacenter local, 1--2ms; entre regiones, 50--200ms.
La decisión de usar replicación síncrona es un tradeoff explícito entre durabilidad y rendimiento.

\begin{lstlisting}[caption={Configurar replicación síncrona}]
-- postgresql.conf en PRIMARY:
synchronous_standby_names = 'FIRST 1 (standby1, standby2)'
-- FIRST 1: el primero de la lista que esté disponible se vuelve síncrono
-- ANY 1: cualquiera de la lista

synchronous_commit = on         -- default; espera flush en standby
-- synchronous_commit = remote_apply  -- más fuerte: espera aplicación

-- Por transacción (tradeoff rendimiento/durabilidad):
BEGIN;
SET LOCAL synchronous_commit = off;  -- esta tx es asíncrona
INSERT INTO audit_log ...;
COMMIT;
\end{lstlisting}

\section[Patroni: HA Automatizado]{\emj{🎯}~Patroni: HA Automatizado}

Patroni es un template de HA para PostgreSQL que usa un DCS (Distributed Configuration Store: etcd, Consul, ZooKeeper) para coordinar elecciones de líder y failover automático.

\begin{teoriabox}{Componentes de Patroni}
\begin{itemize}
  \item \textbf{Patroni agent:} corre en cada nodo de PG. Monitorea el nodo, actualiza el DCS, maneja failover.
  \item \textbf{DCS (etcd):} almacena quién es el líder y el estado del clúster. El ``árbitro'' que evita split-brain.
  \item \textbf{HAProxy/pgBouncer:} balanceador de carga. Dirige escrituras al primary, lecturas a standbys. Patroni expone health endpoints HTTP para que HAProxy detecte el primary.
  \item \textbf{patronictl:} CLI para gestionar el clúster (switchover manual, failover, restart, reinitialización).
\end{itemize}
\end{teoriabox}

El archivo de configuración de Patroni define la identidad del nodo, cómo conectarse al DCS, y los parámetros de PostgreSQL.
Cada nodo del clúster tiene su propio archivo con un \texttt{name} único pero el mismo \texttt{scope} (nombre del clúster).
El \texttt{maximum\_lag\_on\_failover} es especialmente importante: si un standby tiene más lag que este valor en bytes, Patroni no lo promoverá a primary automáticamente, evitando que un standby muy desactualizado se convierta en el nuevo líder.

\begin{lstlisting}[language=, caption={patroni.yml — configuración básica}]
scope: postgres-cluster
namespace: /db/
name: node1

restapi:
  listen: 0.0.0.0:8008
  connect_address: 10.0.0.1:8008

etcd:
  hosts: 10.0.0.10:2379,10.0.0.11:2379,10.0.0.12:2379

bootstrap:
  dcs:
    ttl: 30
    loop_wait: 10
    retry_timeout: 30
    maximum_lag_on_failover: 1048576  # 1MB lag máximo para ser elegible

  initdb:
    - encoding: UTF8
    - data-checksums

postgresql:
  listen: 0.0.0.0:5432
  connect_address: 10.0.0.1:5432
  data_dir: /var/lib/postgresql/data
  authentication:
    replication:
      username: replicator
      password: securepass
    superuser:
      username: postgres
      password: adminpass
\end{lstlisting}

\texttt{patronictl} es la CLI para gestionar el clúster Patroni desde cualquier nodo con acceso al DCS.
Un \textbf{switchover} es un failover planificado: transfiere el liderazgo de forma ordenada, replicando todos los cambios pendientes primero y asegurando mínimo downtime (típicamente $<$1 segundo).
Un \textbf{failover} forzado se usa en emergencias cuando el primary no responde: puede implicar algo de pérdida de datos si había replicación asíncrona con lag.

\begin{lstlisting}[language=, caption={patronictl — operaciones comunes}]
# Ver estado del clúster:
patronictl -c /etc/patroni.yml list

# Switchover manual (planned, sin downtime notable):
patronictl -c /etc/patroni.yml switchover --master node1 --candidate node2

# Failover forzado (emergencia):
patronictl -c /etc/patroni.yml failover --master node1 --force

# Reiniciar nodo:
patronictl -c /etc/patroni.yml restart postgres-cluster node1

# Reinicializar standby (después de desync):
patronictl -c /etc/patroni.yml reinit postgres-cluster node2
\end{lstlisting}

\section[Failover y Split-Brain]{\emj{⚠️}~Failover y Split-Brain}

\begin{warningbox}
\textbf{Split-brain:} dos nodos creen ser el primary simultáneamente. Ambos aceptan escrituras, los datos divergen, y no hay forma automática de reconciliar. Es el peor escenario en HA.

Patroni lo evita con el DCS: el primary tiene un ``lease'' (lock) en etcd. Si pierde conectividad con etcd, \textbf{se degrada a standby} aunque aún pueda hablar con clientes.
El nodo que no puede adquirir el lock no puede ser primary, sin importar su estado de red con los otros nodos de PG.
\end{warningbox}

Patroni expone endpoints HTTP en el puerto 8008 que HAProxy utiliza para health checks.
Esto es fundamental: el balanceador de carga no necesita lógica especial para detectar quién es el primary; simplemente hace un GET cada pocos segundos y enruta el tráfico al nodo que responde 200 OK en \texttt{/primary}.
Cuando ocurre un failover y un standby es promovido, en segundos el health check cambia y HAProxy redirige el tráfico automáticamente.

\begin{lstlisting}[caption={Health endpoint de Patroni para HAProxy}]
-- Patroni expone en :8008:
-- GET /primary   → 200 si es primary, 503 si no
-- GET /replica   → 200 si es standby, 503 si no
-- GET /health    → 200 siempre (para liveness check)

-- haproxy.cfg:
-- backend postgres_primary
--   option httpchk GET /primary
--   server node1 10.0.0.1:5432 check port 8008
--   server node2 10.0.0.2:5432 check port 8008
--   server node3 10.0.0.3:5432 check port 8008
\end{lstlisting}

\begin{entrevistabox}
\begin{enumerate}
  \item ¿Cuál es la diferencia entre RPO y RTO? ¿Cómo los afecta la replicación síncrona vs asíncrona?
  \item Explica cómo Patroni evita split-brain.
  \item ¿Qué es hot\_standby\_feedback y por qué importa en un clúster con queries largas en el standby?
  \item ¿Cuándo usarías synchronous\_commit = remote\_apply en lugar de \texttt{on}?
  \item Un standby tiene 5 minutos de lag. ¿Cuáles son las causas más probables y cómo diagnosticas?
  \item ¿Qué ocurre si etcd no tiene quórum? ¿El clúster sigue funcionando?
\end{enumerate}
\end{entrevistabox}
